% ============================================================
% ESTILOS GENERALES
% Proyecto: Física
% ============================================================


% ------------------------------------------------------------
% TEXTO GENERAL
% ------------------------------------------------------------

\setstretch{1.15}

\setlength{\parindent}{0pt}
\setlength{\parskip}{6pt}

\AtBeginDocument{
    \color{TextoPrincipal}
}


% ------------------------------------------------------------
% LISTAS
% ------------------------------------------------------------

\setlist[itemize]{
    topsep=4pt,
    itemsep=2pt
}

\setlist[enumerate]{
    topsep=4pt,
    itemsep=2pt
}


% ------------------------------------------------------------
% PARTES
% ------------------------------------------------------------

\titleformat{\part}[display]
    {\normalfont\Huge\bfseries\color{ColorTitulos}\centering}
    {\partname\ \thepart}
    {1em}
    {}


% ------------------------------------------------------------
% CAPÍTULOS
% ------------------------------------------------------------

\titleformat{\chapter}
    {\normalfont\huge\bfseries\color{ColorTitulos}}
    {\thechapter.}
    {1em}
    {}


% ------------------------------------------------------------
% SECCIONES
% ------------------------------------------------------------

\titleformat{\section}
    {\normalfont\Large\bfseries\color{ColorTitulos}}
    {\thesection.}
    {0.7em}
    {}

\titleformat{\subsection}
    {\normalfont\large\bfseries\color{ColorTitulos}}
    {\thesubsection.}
    {0.7em}
    {}

\titleformat{\subsubsection}
    {\normalfont\normalsize\bfseries\color{ColorTitulos}}
    {\thesubsubsection.}
    {0.7em}
    {}


% ------------------------------------------------------------
% PROFUNDIDAD DE NUMERACIÓN E ÍNDICE
% ------------------------------------------------------------

\setcounter{secnumdepth}{3}
\setcounter{tocdepth}{3}


% ------------------------------------------------------------
% ENCABEZADOS Y PIES DE PÁGINA
% ------------------------------------------------------------

\pagestyle{fancy}

\fancyhf{}

\fancyhead[L]{
    \nouppercase{\leftmark}
}

\fancyhead[R]{
    \tituloobra
}

\fancyfoot[C]{
    \thepage
}

\renewcommand{\headrulewidth}{0.4pt}
\renewcommand{\footrulewidth}{0pt}

\setlength{\headheight}{15pt}


% ------------------------------------------------------------
% FIGURAS Y TABLAS
% ------------------------------------------------------------

\captionsetup{
    font=small,
    labelfont=bf
}


% ------------------------------------------------------------
% UNIDADES SI
% ------------------------------------------------------------

\sisetup{
    output-decimal-marker={,},
    group-separator={\,},
    group-minimum-digits=4,
    exponent-product=\cdot,
    per-mode=symbol,
    separate-uncertainty=true
}


% ------------------------------------------------------------
% HIPERVÍNCULOS
% ------------------------------------------------------------

\hypersetup{
    colorlinks=true,
    linkcolor=AzulPrincipal,
    citecolor=VerdePrincipal,
    urlcolor=AzulSecundario,
    pdftitle={\tituloobra},
    pdfauthor={\autorlibro},
    pdfsubject={Colección general de Física},
    pdfcreator={LaTeX}
}


% ------------------------------------------------------------
% ESTILOS BASE DE TCOLORBOX
% ------------------------------------------------------------
%
% IMPORTANTE:
% No introducir líneas en blanco dentro de \tcbset{...}.
% pgfkeys puede interpretar una línea en blanco como \par
% y producir errores de análisis de claves.
% ------------------------------------------------------------

\tcbset{
    estiloBase/.style={
        enhanced,
        breakable,
        boxrule=0.7pt,
        arc=2mm,
        left=3mm,
        right=3mm,
        top=2mm,
        bottom=2mm,
        before skip=10pt,
        after skip=10pt,
        fonttitle=\bfseries
    },
    estiloRespuesta/.style={
        enhanced,
        breakable,
        colback=GrisFondo,
        colframe=GrisPrincipal,
        boxrule=0.6pt,
        arc=1.5mm,
        left=3mm,
        right=3mm,
        top=2mm,
        bottom=2mm,
        before skip=6pt,
        after skip=6pt,
        title=Respuesta,
        fonttitle=\bfseries
    },
    estiloSolucion/.style={
        enhanced,
        breakable,
        colback=VerdeFondo,
        colframe=VerdePrincipal,
        boxrule=0.6pt,
        arc=1.5mm,
        left=3mm,
        right=3mm,
        top=2mm,
        bottom=2mm,
        before skip=6pt,
        after skip=10pt,
        title=Solución desarrollada,
        fonttitle=\bfseries
    }
}